\section{Reproducing the paper}
\label{sec:reproduce}

\begin{shell}
build/bin/experiment all\\
build/bin/baselines results/baselines.csv results/baselines\_trace.csv\\
for p in slow constraints domain; do build/bin/baselines b\_\$p.csv t\_\$p.csv -{}-platform \$p; done\\
(cat b\_slow.csv; tail -n +2 b\_constraints.csv; tail -n +2 b\_domain.csv) > results/baselines\_platforms.csv\\
python3 exact/gap.py build/bin/export\_instance -{}-sizes 20,50,100,200,300 -{}-seeds 10\\
python3 exact/gap.py build/bin/export\_instance -{}-sizes 20,50,100,200,300 -{}-seeds 10 \textbackslash\\
\quad -{}-configs domain,hybrid -{}-out results/exact\_architectures.csv\\
sh tests/netsim/run\_netsim.sh build/bin/net\_export tests/netsim/netsim.py 20000 2\\
sh tests/syssim/experiment.sh build/bin/mapping results/syssim.csv 5 4\\
cd paper \&\& make figures \&\& make
\end{shell}
The first command writes the CSV files of the 24 experiments; the next
three the metaheuristic comparison on the default platform and on the
three harder ones (about twenty minutes each on two cores); the
following two the exact results on the zonal platforms (about fifteen
minutes) and on the domain and hybrid ones (about three minutes); the
next two print the simulation checks and write \cmd{results/syssim.csv}; the last
draws the figures, writes the generated tables
(\cmd{paper/make\_tables.py}) and builds \cmd{paper/gauthier2026-system-level-mapping-automotive.pdf}. The results
are deterministic: the same commands give the same files on any machine
and with any number of threads, CPU times aside.

\paragraph{The use cases and this manual.}
\begin{shell}
cd manual \&\& make listings \&\& make
\end{shell}
regenerates every output shown in \cref{sec:output,sec:usecases} from the
files of \cmd{examples/}, then the manual.
